\section{Conclusion}

\subsection{Answers to the Research Questions}

\textbf{RQ1: deterministic and safe transformation.} Permission summaries are accepted through a conservative evidence contract, rejected when identifiers, numbers, sources, windows, timestamps, or contexts are unsafe, and evaluated through a deterministic rule-pack engine. Temporal isolation addresses replay and overlap within the limits of aggregate windows.

\textbf{RQ2: calibrated communication.} Each finding preserves provenance, temporal scope, rule-pack identity, rationale, and missing-context questions. The interface distinguishes needs-review, evidence-gap, and no-technical-concern states, and repeatedly states that technical output is not a legal verdict. Synthetic demonstrations remain labelled throughout the data path.

\textbf{RQ3: regional extensibility.} Jurisdiction-dependent material is loaded from registered, versioned packs behind a stable contract. The EU GDPR pack is the evaluated legal objective; the Research Baseline demonstrates software switching without claiming another jurisdiction. New regional packs require their own legal and empirical validation.

\textbf{RQ4: supported claims.} Compilation, lint, deterministic and adversarial suites, release APK/AAB assembly, merged-manifest inspection, physical installation, cold launch, and rendered workflows passed for the recorded candidate. These results support an initial store-submission engineering candidate. They do not prove legal compliance, unrestricted Android visibility, long-run reliability, broad usability, accessibility conformance, production signing, or store acceptance.

\subsection{Contributions}

The thesis contributes:

\begin{enumerate}
    \item an evidence-bounded Android privacy-review architecture;
    \item a fail-closed, temporally isolated decision pipeline with explicit provenance;
    \item a replaceable regulation-pack contract that separates legal mappings from product infrastructure;
    \item a local-first, three-destination interface designed for calibrated trust;
    \item a reproducible verification record that separates logical, package, physical-device, legal, and deployment claims.
\end{enumerate}

\subsection{Final Product Position}

Privacy Lens 1.1.0 was judged visually coherent, trustworthy within its disclosed boundaries, and usable enough to invite continued use in a controlled initial-submission context. The final release artefacts use package \texttt{com.zhihengzhang.privacylens}, target SDK 36, and request no sensitive runtime permission. They are accompanied by a privacy-policy draft, user guide, product audit, store-readiness checklist, checksums, and physical-device screenshots.

The candidate is not yet a publicly publishable production release. Production signing, stable public policy and support endpoints, console declarations, internal-track evaluation, and store review remain external owner actions. This distinction is part of the result rather than an administrative footnote: trustworthy privacy software must disclose where its evidence ends.

\subsection{Future Work}

Future work should validate the GDPR pack independently, add a genuinely reviewed regional pack, improve authorised event-level acquisition, test multiple Android implementations over long durations, and conduct accessibility and participant studies. These activities should retain the Revision 9 claim-to-evidence discipline. Additional features are valuable only when their authority, provenance, and validation remain visible.

\subsection{Closing Statement}

The central finding is that privacy accountability on Android is strengthened not by presenting the strongest possible conclusion, but by presenting the strongest conclusion the available evidence actually supports. Privacy Lens operationalises that principle in code, interface, evaluation, and thesis structure.
